\documentclass[11pt]{article}
% Uses the canonical latex-docs semantic style in-repository. The local
% lightweight fallback is for independent compilation of this distribution.
\IfFileExists{technical-implementation.sty}{\usepackage{technical-implementation}}{\usepackage{../support/aaoa-standalone}}
\usepackage{mathtools}
\setDocTitle{Generating Functions and Counting Workbook -- Blank}
\setDocSubtitle{Printable, writable assessment copy}
\setDocVersion{1.0}
\setDocStatus{Educational reference / unapproved working material}
\setDocOwner{Algorithm Engineering}
\setDocAuthor{Original companion material}
\setDocSummary{Printable workbook with specified problems and empty answer fields; no inserted answers or bracketed prompts.}
\providecommand{\Expect}{\mathbb{E}}
\providecommand{\Prob}{\mathbb{P}}
\providecommand{\Var}{\operatorname{Var}}
\providecommand{\bigoh}{\mathcal{O}}
\begin{document}
\makedoctitle
\section{Participant and session record}
\begin{tabularx}{\linewidth}{@{}lX@{}}Date: & \hrulefill\\ Participant: & \hrulefill\\ Source chapter: & \hrulefill\\\end{tabularx}
\section{Problem 1: $a_0=2$, $a_n=2a_{n-1}+3$}
Derive OGF and extract coefficients. Test first six values.\par\medskip
\textbf{Assumptions and definitions}\par\vspace{0.7cm}\hrule\vspace{0.35cm}
\textbf{Derivation}\par\vspace{2.1cm}\hrule\vspace{0.3cm}
\textbf{Boundary and numerical checks}\par\vspace{1.1cm}\hrule\vspace{0.35cm}
\textbf{Limitations}\par\vspace{0.9cm}\hrule\vspace{0.35cm}
\section{Problem 2: $A(z)=(1-z)^{-4}$}
Derive a finite coefficient formula and the first five terms.\par\medskip
\textbf{Assumptions and definitions}\par\vspace{0.7cm}\hrule\vspace{0.35cm}
\textbf{Derivation}\par\vspace{2.1cm}\hrule\vspace{0.3cm}
\textbf{Boundary and numerical checks}\par\vspace{1.1cm}\hrule\vspace{0.35cm}
\textbf{Limitations}\par\vspace{0.9cm}\hrule\vspace{0.35cm}
\section{Problem 3: An ordered binary tree is empty or root-plus-two-ordered-subtrees}
Write a combinatorial class equation, OGF, exact coefficients for $n=0,\ldots,5$, and the limiting scale.\par\medskip
\textbf{Assumptions and definitions}\par\vspace{0.7cm}\hrule\vspace{0.35cm}
\textbf{Derivation}\par\vspace{2.1cm}\hrule\vspace{0.3cm}
\textbf{Boundary and numerical checks}\par\vspace{1.1cm}\hrule\vspace{0.35cm}
\textbf{Limitations}\par\vspace{0.9cm}\hrule\vspace{0.35cm}
\section{Problem 4: Four independent Bernoulli$(p)$ trials}
Construct a PGF and use derivatives to obtain mean and variance.\par\medskip
\textbf{Assumptions and definitions}\par\vspace{0.7cm}\hrule\vspace{0.35cm}
\textbf{Derivation}\par\vspace{2.1cm}\hrule\vspace{0.3cm}
\textbf{Boundary and numerical checks}\par\vspace{1.1cm}\hrule\vspace{0.35cm}
\textbf{Limitations}\par\vspace{0.9cm}\hrule\vspace{0.35cm}
\section{Problem 5: $F(z,u)$ marks tree size and a parameter}
Explain which coefficient ratio yields the expected parameter and why $u=1$ is necessary.\par\medskip
\textbf{Assumptions and definitions}\par\vspace{0.7cm}\hrule\vspace{0.35cm}
\textbf{Derivation}\par\vspace{2.1cm}\hrule\vspace{0.3cm}
\textbf{Boundary and numerical checks}\par\vspace{1.1cm}\hrule\vspace{0.35cm}
\textbf{Limitations}\par\vspace{0.9cm}\hrule\vspace{0.35cm}
\clearpage
\section{Final quality check}
\begin{tabularx}{\linewidth}{@{}lX@{}}Checked by: & \hrulefill\\ Date: & \hrulefill\\ Additional evidence: & \hrulefill\\\end{tabularx}
\vspace{3cm}\hrule

\end{document}
